% TODO: Finish up this chapter
\section{Basics Of Computer Networks}

\bd[Computer Network]
A \textbf{computer network} (or more simply \textbf{network}) is a system that allows multiple connected devices to
communicate with one another, exchange data, and share information.
\ed

Computer networks support many applications and services, such as access to the World Wide Web, shared use of
application and storage servers, and use of email and instant messaging applications. \v

Each piece of device connected to a network is called a ``node''.

\bd[Network Node]
\textbf{Network nodes} (or more simply \textbf{nodes}) are the physical pieces that are attached to a network, and are
capable of creating, receiving, or transmitting information over a communication channel.
\ed

The nodes of a network may include personal computers, servers, networking hardware and many more. \v

Nodes in a network are connected by communication links. These communication links can be established over cables or
wireless media.

\bd[Communication Link]
A \textbf{communication link} is a communication channel that connects two or more nodes in a network.
\ed

These communication links can be established over cables or wireless media. Based on if networks use cables or not,
we can divide networks into two categories: ``wired'' and ``wireless''.

\bd[Wired Network]
A \textbf{wired network} is a network that uses cables to connect the nodes.
\ed

\bd[Wireless Network]
A \textbf{wireless network} is a network that uses wireless data connections between nodes.
\ed

\bd[Network Topology]
The arrangement of the nodes and their links in a network is called \textbf{network topology}.
\ed

\subsection{Types Of Computer Networks}

There are several types of networks depending on their range. Here follow some of the most important ones.
\bit
\item \textbf{Body Area Network (BAN)}: A Body Area Network (BAN) (or in the case of wireless a Wireless Body Area
Network (WBAN)) is a wired/wireless network of wearable computing devices.
\item \textbf{Personal Area Network (PAN)}: A Personal Area Network (PAN) is a network for interconnecting electronic
devices within an individual person's workspace. A PAN provides data transmission among devices such as computers,
smartphones, tablets and personal digital assistants.
\item \textbf{Small Office/Home Office (SOHO)}: Small Office/Home Office (SOHO) networks are small networks typically
consisted of less than 10 computers.
\item \textbf{Local Area Network (LAN)}: A Local Area Network (LAN) (or in the case of wireless a Wireless Local Area
Network (WLAN)) is a network that interconnects computers within a limited area such as a residence, school,
laboratory, university campus or office building.
\item \textbf{Campus Area Network (CAN)}: A Campus Area Network (CAN) is a network made up of an interconnection of
local area networks (LANs) within a limited geographical area.
\item \textbf{Metropolitan Area Network (MAN)}: A Metropolitan Area Network (MAN) is a computer network that
interconnects users with computer resources in a geographic region of the size of a metropolitan area.
\item \textbf{Wide Area Network (WAN)}: A Wide Area Network (WAN) is a telecommunications network that extends over a
large geographic area for the primary purpose of computer networking.
\eit

Here is a graph summarizing all the type of networks ordered by their range. \v

\fig{net}{0.68}

Finally, there is the concept of a ``network service provider''.

\bd[Network Service Provider (NSP)]
A \textbf{Network Service Provider (NSP)} is a business or organization that sells bandwidth or network access by
providing direct Internet backbone access to internet service providers and usually access to its network access
points.
\ed

\subsection{Client–Server Model}

Nodes are divided into two categories: ``servers'' and ``clients''.

\bd[Server]
A \textbf{server} is a node that provides functionality for other nodes in a network.
\ed

These functionalities are often called ``services''.

\bd[Service]
A \textbf{service} is a piece of functionality a server provides.
\ed

\bd[Client]
A \textbf{client} is a node that accesses a service made available by a server in a computer network.
\ed

Servers can provide various services such as sharing data or resources among multiple clients, or performing
computation for a client. A single server can serve multiple clients, and a single client can use multiple servers. A
client process may run on the same device or may connect over a network to a server on a different device. \v

Servers and clients are integrated in the so called ``client-server'' model.

\bd[Client–Server Model]
The \textbf{client–server model} is a distributed application structure that partitions tasks or workloads between
the providers of a resource or service, called servers, and service requesters, called clients.
\ed

\subsection{Protocols \& Standards}

All nodes in a network must be able to communicate with each other. In order for this to be able they must all
``speak'' the same ``language''. In networks, these ``languages'' are called ``protocols''.

\bd[Protocol]
A \textbf{protocol} is a system of rules and instructions that defines the syntax, semantics and synchronization of
communication (and possible error recovery methods) among the entities of a network and thus allows the transmission
of information.
\ed

\bd[Protocol Data Unit (PDU)]
A \textbf{Protocol Data Unit (PDU)} is a single unit of information transmitted among peer entities of a network.
\ed

A PDU is composed of protocol-specific control information and user data. In the layered architectures of
communication protocol stacks, each layer implements protocols tailored to the specific type or mode of data exchange.

\bd[Standard]
A \textbf{standard} is a formalized protocol accepted by most of the parties that implement it.
\ed

There are many protocols and standards depending on the task in hand.

\section{Open Systems Interconnection (OSI)}

\bd[Open Systems Interconnection (OSI)]
The \textbf{Open Systems Interconnection Model (OSI)} is a conceptual framework used to describe the functions of a
networking system. The OSI model characterizes computing functions into a universal set of rules and requirements in
order to support interoperability between different products and software.
\ed

OSI developed in the late 1970s to support the emergence of the diverse networks that were competing for application
in the large national networking efforts in the world, is a conceptual model that characterizes and standardises the
communication functions of a network without regard to its underlying internal structure and technology. Its goal is
the interoperability of diverse communication systems with standard communication protocols. \v

The model partitions the flow of data in a network into seven abstraction layers, from the physical implementation of
transmitting bits across a communications medium to the highest-level representation of data of a distributed
application. Each intermediate layer serves a class of functionality to the layer above it and is served by the layer
below it. \v

Classes of functionality are realized in software by standardised communication protocols which enable an entity in
one host to interact with a corresponding entity at the same layer in another host. OSI abstractly describe the
functionality provided to an ($N$) layer by an ($N-1$) layer, where $N$ is one of the seven layers of protocols
operating in the local host. At each level $N$, two entities at the communicating devices (layer $N$ peers) exchange
PDUs by means of a layer $N$ protocol. \v

Here are the 7 layers of OSI model in a diagram.

\fig{net1}{0.45}

\subsection{Layer 1: Physical Layer}

\bd[Physical Layer]
The \textbf{physical layer} is the first layer of the OSI model and defines the means of transmitting a stream of raw
bits (physical layer's PDU) over a physical data link connecting network nodes.
\ed

The physical layer is the first and lowest layer in the OSI model, and it is the layer most closely associated with
the physical connection between devices, providing an electrical, mechanical, and procedural interface to the
transmission medium. The physical layer is responsible for the transmission and reception of unstructured streams of
raw bits between a device and a physical transmission medium. It then converts the digital bits into electrical,
radio, or optical signals. It consists of the electronic circuit transmission technologies of a network. It is a
fundamental layer underlying the higher level functions in a network, and can be implemented through a great number
of different hardware technologies with widely varying characteristics. Layer specifications define characteristics
such as voltage levels, the timing of voltage changes, physical data rates, maximum transmission distances,
modulation scheme, channel access method and physical connectors. This includes the layout of pins, voltages, line
impedance, cable specifications, signal timing and frequency for wireless devices. Bit rate control is done at the
physical layer and may define transmission mode as simplex, half duplex, and full duplex. The components of a
physical layer can be described in terms of a network topology. \v

Here follows a list of the most common protocols and standards of the physical layer.
\bit
\item \textbf{Integrated Services Digital Network (ISDN)} is a set of communication standards for simultaneous
digital transmission of voice, video, data, and other network services over the digitalised circuits of the public
switched telephone network. Work on the standard began in 1980 at Bell Labs and was formally standardised in 1988. By
the time the standard was released, newer networking systems with much greater speeds were available, and ISDN saw
relatively little uptake in the wider market. ISDN has largely been replaced with digital subscriber line (DSL)
systems of much higher performance.
\item \textbf{Digital Subscriber Line (DSL)} is a family of technologies that are used to transmit digital data over
telephone lines. DSL service can be delivered simultaneously with wired telephone service on the same telephone line
since DSL uses higher frequency bands for data. The bit rate of consumer DSL services typically ranges from 256
kbit/s to over 100 Mbit/s in the direction to the customer (downstream), depending on DSL technology, line
conditions, and service-level implementation.
\item \textbf{Asymmetric Digital Subscriber Line (ADSL)} the most commonly installed DSL technology, for Internet
access. In ADSL, the data throughput in the upstream direction (the direction to the service provider) is lower,
hence the designation of asymmetric service.
\item \textbf{IEEE \footnote{IEEE 802 is a family of IEEE (Institute of Electrical and Electronics Engineers) protocols
and standards for local area networks (LAN), personal area network (PAN), and metropolitan area networks (MAN). The
number 802 has no significance, it was simply the next number in the sequence that the IEEE used for standards projects.
The services and protocols specified in IEEE 802 map to the lower two layers(physical and data link) of OSI. The IEEE
802 family of standards has twelve members, numbered 802.1 through 802.12} 802.3} or \textbf{Ethernet} where nodes in
the LAN connect to a central router, hub, switch, or gateway which acts as a connecting point between devices on the
network through Unshielded Twisted Pair (UTP) cables. In this protocol, when a node on the network wants to send data to
another node, it senses the carrier, which is the main wire connecting the devices. If it is free, meaning no one is
sending anything, it sends the data packet on the network, and the other devices check the packet to see whether they
are the recipient. The recipient consumes the packet. If there is a packet on the highway, the device that wants to
send holds back for some thousandths of a second to try again until it can send.
\item \textbf{IEEE 802.11} or \textbf{Wireless LAN (WLAN)} tandards are the most widely used computer networks in the
world. These are commonly called Wi-Fi, which is a trademark belonging to the Wi-Fi Alliance.
\item \textbf{Universal Serial Bus (USB)} is an industry standard that establishes specifications for cables,
connectors and protocols for connection, communication and power supply (interfacing) between computers, peripherals
and other computers. Released in 1996, there have been four generations of USB specifications: USB 1.x, USB 2.0, USB
3.x, and USB4. A broad variety of USB hardware exists, including fourteen different connectors, of which USB-C, as of
now, is the most recent and best one.
\eit

\subsection{Layer 2: Data Link Layer}

\bd[Data Link Layer]
The \textbf{data link layer} is the second layer of the OSI model, and it is the protocol layer that transfers data
in the form of frames (data link layer's PDU) between nodes on a network segment across the physical layer.
\ed

The data link layer is the second layer and is the protocol layer that provides the functional and procedural means
to transfer data between directly connected nodes on a network segment across the physical layer. It also detects and
possibly corrects errors that may occur in the physical layer. The data link layer is concerned with local delivery
of frames between nodes on the same level of the network. Data-link frames, as these protocol data units are called,
do not cross the boundaries of a LAN. \v

IEEE 802 divides the data link layer into two sub-layers.

\bd[Medium Access Control (MAC) Sublayer]
The \textbf{Medium Access Control (MAC) sublayer} is the sublayer that controls the hardware which is responsible for
controlling how devices in a network gain access to a medium and permission to transmit data.
\ed

MAC sublayer provides flow control and multiplexing for the transmission medium. When sending data to another device
on the network, the MAC sublayer encapsulates higher-level frames into frames appropriate for the transmission
medium, adds a frame check sequence to identify transmission errors, and then forwards the data to the physical layer
as soon as the appropriate channel access method permits it. Additionally, the MAC is also responsible for
compensating for collisions by initiating retransmission if a jam signal is detected. When receiving data from the
physical layer, the MAC block ensures data integrity by verifying the sender's frame check sequences, and strips off
the sender's preamble and padding before passing the data up to the higher layers.

\bd[Logical Link Control (LLC) Sublayer]
The \textbf{Logical link control (LLC) sublayer} is the upper sublayer of the data link layer, which acts as an
interface between the MAC sublayer and the network layer.
\ed

LLC sublayer is responsible for identifying and encapsulating network layer protocols, and controls error checking
and frame synchronization. It provides multiplexing mechanisms that make it possible for several network protocols to
coexist within a multipoint network and to be transported over the same network medium. It can also provide
node-to-node flow control and error management.

\subsection{Layer 3: Network Layer}

\subsection{Layer 4: Transport Layer}

\subsection{Layer 5: Session Layer}

\subsection{Layer 6: Presentation Layer}

\subsection{Layer 7: Application Layer}

\newpage

\section{Network Devices (Hardware)}

\bd[Network Devices (Hardware)]
\textbf{Network devices} (or \textbf{hardware}) are physical devices that are required for communication and interaction
between hardware on a computer network.
\ed

Network devices or hardware are the physical representation of nodes. There are many network devices. Let's take a
look at some of the most important ones.

\bd[Cables]
\textbf{Cables} are networking devices used to connect one network device to other network devices.
\ed

Cables can be made either of copper or fiber. Copper cables use electric signals, are cheaper and used for short
distances. Due to their nature they can be affected by electromagnetic interference. Fiber cables are made of glass,
use light signals, are more expensive and used for longer distances. Due to their nature they are not affected by
electromagnetic interference.

\bd[Hub]
A \textbf{hub} is the most basic hardware that connects multiple network devices together and acts as a repeater in
that it amplifies signals that deteriorate after traveling long distances over connecting cables.
\ed

\fig{hub}{0.2}

A hub is the simplest in the family of network connecting devices because it connects LAN components with identical
protocols. A hub can be used with both digital and analog data, provided its settings have been configured to prepare
for the formatting of the incoming data. \v

Hubs do not perform packet filtering or addressing functions; they just send data packets to all connected devices.
Hubs operate at the Physical layer of the Open Systems Interconnection (OSI) model. (More on that later) \v

Hubs are quite old technology, and they are not used anymore. Their main problem is something called collision domain.

\bd[Collision Domain]
A \textbf{collision domain} is a network segment connected by a shared medium or through repeaters where simultaneous
data transmissions collide with one another.
\ed

A network collision occurs when more than one device attempts to send a packet on a network segment at the same time.
Members of a collision domain may be involved in collisions with one another. Devices outside the collision domain do
not have collisions with those inside.

\bd[Bridge]
A \textbf{bridge} is a computer networking device that creates a single,aggregate network from multiple communication
networks or network segments.
\ed

\fig{bridge}{0.5}

Bridges are used to connect two or more hosts or network segments together.The basic role of bridges in network
architecture is storing and forwarding frames between the different segments that the bridge connects. They use
hardware Media Access Control (MAC) addresses for transferring frames. By looking at the MAC address of the devices
connected to each segment, bridges can forward the data or block it from crossing. Bridges can also be used to
connect two physical LANs into a larger logical LAN. \v

Bridges work only at the Physical and Data Link layers of the OSI model. Bridges are used to divide larger networks
into smaller sections by sitting between two physical network segments and managing the flow of data between the two. \v

Bridges are like hubs in many respects, including the fact that they connect LAN components with identical protocols.
However, bridges filter incoming data packets, known as frames, for addresses before they are forwarded. As it
filters the data packets, the bridge makes no modifications to the format or content of the incoming data. The bridge
filters and forwards frames on the network with the help of a dynamic bridge table. The bridge table, which is
initially empty, maintains the LAN addresses for each computer in the LAN and the addresses of each bridge interface
that connects the LAN to other LANs. Bridges, like hubs, can be either simple or multiple port. \v

Bridges have mostly fallen out of favor in recent years and have been replaced by switches, which offer more
functionality. In fact, switches are sometimes referred to as ``multiport bridges'' because of how they operate.

\bd[Switch]
A \textbf{switch} is hardware that connects devices on a computer network by using packet switching to receive and
forward data to the destination device.
\ed

\fig{switch}{0.65}

Switches generally have a more intelligent role than hubs. A switch is a multiport device that improves network
efficiency. The switch maintains limited routing information about nodes in the internal network, and it allows
connections to systems like hubs. Strands of LANs are usually connected using switches. Generally, switches can read
the hardware addresses of incoming packets to transmit them to the appropriate destination. \v

Using switches improves network efficiency over hubs because of the virtual circuit capability. Switches also improve
network security because the virtual circuits are more difficult to examine with network monitors. You can think of a
switch as a device that has some of the best capabilities of hubs. A switch can work at either the Data Link layer or
the Network layer of the OSI model. A multilayer switch is one that can operate at both layers, which means that it
can operate as both a switch and a router which we will define right next. \v

Hubs, bridges and switches are the primary devices used to connect network devices on a single network, usually a LAN\@.
However, we often want to send or receive data between independent layers. That's where routers are coming in use.

\bd[Router]
A \textbf{router} is a networking device that forwards data packets between computer networks.
\ed

\fig{router}{0.34}

Routers help transmit packets to their destinations by charting a path through the sea of interconnected networking
devices using different network topologies. Routers are intelligent devices, and they store information about the
networks they're connected to. Most routers can be configured to operate as packet-filtering firewalls and use access
control lists. Routers are also used to translate from LAN framing to WAN framing. This is needed because LANs and
WANs use different network protocols. \v

Router are also used to divide internal networks into two or more subnetworks. Routers can also be connected
internally to other routers, creating zones that operate independently. Routers establish communication by
maintaining tables about destinations and local connections. A router contains information about the systems
connected to it and where to send requests if the destination isn't known. \v

Routers are your first line of defense, and they must be configured to pass only traffic that is authorized by
network administrators. The routes themselves can be configured as static or dynamic. If they are static, they can
only be configured manually and stay that way until changed. If they are dynamic, they learn of other routers around
them and use information about those routers to build their routing tables. Routers work at the Network layer of the
OSI model.

\bd[Gateway]
A \textbf{gateway} is a piece of hardware used in that allows data to flow from one discrete network to another.
\ed

\fig{gateway}{0.2}

Gateways are distinct from routers or switches in that they communicate using more than one protocol to connect
multiple networks and can operate at any of the seven layers of OSI, although they normally work at the Transport and
Session layers. At the Transport layer and above, there are numerous protocols and standards from different vendors;
gateways are used to deal with them. Gateways provide translation between networking technologies such as OSI and
TCP/IP. Because of this, gateways connect two or more autonomous networks, each with its own routing algorithms,
protocols, topology, domain name service, and network administration procedures and policies. \v

Gateways perform all the functions of routers and more. In fact, a router with added translation functionality is a
gateway. The function that does the translation between different network technologies is called a protocol converter. \v

The term gateway can also loosely refer to a computer or computer program configured to perform the tasks of a
gateway, such as a default gateway or router, and in the case of HTTP, gateway is also often used as a synonym for
reverse proxy.

\bd[Modem]
A \textbf{modulator-demodulator} (or simply \textbf{modem}), is a hardware device that converts data from a digital
format into a format suitable for an analog device such as computers.
\ed

\fig{modem}{0.24}

A modem transmits data by modulating one or more carrier wave signals to encode digital information, while the
receiver demodulates the signal to recreate the original digital information. The goal is to produce an electrical
signal that can be transmitted easily and decoded reliably. Modems can be used with almost any means of transmitting
analog signals, from light-emitting diodes to radio. Modems work on both the Physical and Data Link layers. \v

The Institute of Electrical and Electronics Engineers (IEEE) is a professional association for electronic and
electrical engineering formed in 1963 with its corporate office in New York City. IEEE 802 is a family of IEEE's
standards for local area networks (LAN), personal area network (PAN), and metropolitan area networks (MAN). (The
number 802 has no significance, it was simply the next number in the sequence that the IEEE used for standards
projects). The IEEE 802 family of standards has had twenty-four members, numbered 802.1 through 802.24, however, not
all of these are currently active. Two of the most widely used, and still active, of these standards are the IEEE 802.3
for the Ethernet family, and the IEEE 802.11 for Wireless LAN (WLAN) \& Mesh (Wi-Fi certification). \v

IEEE 802.3 is a working group and a collection of Institute of Electrical and Electronics Engineers (IEEE) standards
produced by the working group defining the physical layer and data link layer's media access control (MAC) of wired
Ethernet. This is generally a local area network (LAN) technology with some wide area network (WAN) applications.
Physical connections are made between nodes and/or infrastructure devices (hubs, switches, routers) by various types
of copper or fiber cable. \v

The services and protocols specified in IEEE 802 map to the lower two layers (data link and physical) of the
seven-layer Open Systems Interconnection (OSI) networking reference model. More specifically IEEE 802 divides the OSI
data link layer into two sub-layers: logical link control (LLC) and medium access control (MAC).